\section{Economic significance}\label{sec:strategies}

The pre-registered strategies turn each reading into a daily long-short book within size tiers, hedged against
the market with SPY, and ask whether the readings are worth trading. The \emph{convergence strategy} buys the
events whose price lies furthest below the model's value (the top 20\% of the gap $U$) and sells those furthest
above it, after the first half hour (D1) or after the close (D2). Three more families rank the pre-open reading
(entering at the open), and the post-reaction judgment $J$, signed by the move, at 10:00 and after the close.
The tables label the books by their entry point: S-D0 for the pre-open reading, S-D1 and S-D2 for the judgment at
10:00 and after the close, and S-CONV for the convergence strategy; an RG book trades only the events the model
calls an under- or over-reaction, and a gated book only its most confident readings. Costs are charged at 5, 10 and 20 basis points per side on both legs and on the hedge. Each book is compared
with 100 books built on random signals over the same events, and its Sharpe ratio is deflated for the number of
trials behind it \citep{bailey2014}. We deflate with the \res{ntrials-strat} trials the ledger held when the
strategies were evaluated: every selection of filters, encodings, models and baselines, and every strategy
variant up to then. The exploratory study of Section~\ref{sec:exploratory} later raised the ledger to
\res{ntrials-total} trials, which would only raise the bar (Appendix~\ref{app:dsr} gives the formula).

No pre-registered strategy is worth trading after costs (Table~\ref{tab:strategies} and
Figure~\ref{fig:strategies}). Gross of costs, no book reaches the 95th percentile of its random-signal placebo,
and Laya-Fin's post-reaction book after the close falls below the 5th (percentile \res{sd2j-pct}), as expected
of judgments that carry no positive information about the next week. Net of 10 basis points every main book
loses except the tiny book on the model's under- and over-reaction calls, which holds one stock a side on a
handful of days. The convergence strategy earns a gross Sharpe ratio of \res{sconv-gross} for Laya-Fin and
\res{sconv-net} net of costs. It is, in effect, a price-only gap fade: because the model's value is small next
to the day's move, ranking by the gap $U = V - \Glvl$ buys the day's losers and sells its winners. The two books
share \res{sconv-shared} of their positions, \res{sconv-same-side} of them on the same side, and their daily
returns correlate \res{sconv-corr}. The highest gross Sharpe ratio, \res{sd0-gated-gross} for the pre-open book
restricted to Laya-Fin's most confident readings and held to the next close, reaches percentile
\res{sd0-gated-pct} of its placebo and breaks even at \res{sd0-gated-be} basis points per side. The one-day books turn
over about \res{sd0-turnover} times their capital a day, so costs decide their sign. With \res{ntrials-strat}
trials the deflation benchmark is a daily Sharpe ratio of \res{dsr-sr0}, and the largest deflated-Sharpe
probability of any test book is \res{dsr-max}. Short-sale borrow costs, which would lower every net figure, are
not modelled.

\input{\tabdir/tab6_strategies.tex}

\begin{figure}[tbp]
\centering
\includegraphics{\figdir/fig6_strategies.pdf}
\caption{\textbf{No pre-registered strategy survives costs: net of 10 basis points every main book loses, and
the convergence strategy is effectively a price-only gap fade (\res{sconv-shared} of positions shared, daily
correlation \res{sconv-corr}).} Test period. (a, b) Annualized Sharpe ratio of each main book, gross and net of
10 basis points per side on both legs and the SPY hedge, at the exit of its hypothesis' horizon; the gray box is
the 5th to 95th percentile of 100 random-signal books on the same events; (b) has its own scale. (c) The
convergence strategy's cumulative net return (Laya-Fin and d1-ultra) against the price-only gap fade, with the
placebo's 5th to 95th percentile band. Every book is in the replication package. d1-ultra is
model-specific.}\label{fig:strategies}
\end{figure}
